% Experiment 2 (main table): backend/evals/paper/exp2_collusion.py main <regimes>  (30 seeds)
\emph{Regimes.} We compare ten detectors in five behavioural regimes, 30 seeds each.
In the \emph{routine} regime friends meet only through the router. The \emph{busy}
regime has twice as many runners online and friends four times likelier to accept
each other's offers, so the router itself puts honest friends together. In the
\emph{favours} regime every honest group member also posts six errands per semester
meant for a particular friend, who takes them from the feed; this is deliberate
in-group routing with nothing dishonest about it. In the \emph{coordinated} regime the
rings are careful: before each farming errand the partner comes online and walks
over, so the router offers the errand to the partner first. The \emph{hard} regime
combines honest favours with coordinated rings.

\emph{Detectors.} Four single-signal baselines use structure or money at the
production knees. A cycle baseline applies the ring search's leg rule and
strongly-connected-component search to any pair of students, without the friendship
test. Two graph-mining baselines come from the fraud literature. The first is
FRAUDAR-style greedy peeling \cite{hooi2016fraudar,charikar2000dense} on the
bipartite payer--payee graph, with edges weighted by the number of errands and the
logarithmic column weighting FRAUDAR uses against camouflage. The second is
label-propagation communities \cite{raghavan2007lpa} on the friendship graph, ranked
by the share of their members' value that stays inside. The last three rows are the
ring search and the ring search followed by each version of the policy test.

\begin{table*}[!t]
\caption{Ring Detection on a Simulated 1,200-Student Campus. Precision (P), Recall
(R) and False-Positive Rate (FPR, \%), Mean of 30 Seeds}
\label{tab:rings}
\centering
\scriptsize
\setlength{\tabcolsep}{2.4pt}
\begin{tabular}{@{}lccccccccccccccc@{}}
\toprule
& \multicolumn{3}{c}{Routine} & \multicolumn{3}{c}{Busy} & \multicolumn{3}{c}{Favours}
& \multicolumn{3}{c}{Coordinated} & \multicolumn{3}{c}{Hard} \\
\cmidrule(lr){2-4}\cmidrule(lr){5-7}\cmidrule(lr){8-10}\cmidrule(lr){11-13}\cmidrule(l){14-16}
Detector & P & R & FPR & P & R & FPR & P & R & FPR & P & R & FPR & P & R & FPR \\
\midrule
Closure, circulation, or both & n/a & 0.00 & 0.00 & n/a & 0.00 & 0.03 & n/a & 0.00 & 0.00 & n/a & 0.00 & 0.00 & n/a & 0.00 & 0.00 \\
Clustering coefficient $>0.55$ & 0.11 & 0.03 & 0.44 & 0.11 & 0.03 & 0.44 & 0.11 & 0.03 & 0.44 & 0.11 & 0.03 & 0.44 & 0.11 & 0.03 & 0.44 \\
Payment cycles, any pair & 0.63 & 0.95 & 1.34 & 0.38 & 0.96 & 3.73 & 0.50 & 0.94 & 2.23 & 0.69 & 0.99 & 1.12 & 0.48 & 0.99 & 2.48 \\
FRAUDAR-style peeling$^{\dagger}$ & 0.36 & 0.36 & 1.41 & 0.28 & 0.28 & 1.59 & 0.30 & 0.30 & 1.55 & 0.48 & 0.48 & 1.16 & 0.43 & 0.43 & 1.27 \\
Label propagation + money$^{\dagger}$ & 0.03 & 0.06 & 4.18 & 0.02 & 0.03 & 4.17 & 0.02 & 0.04 & 4.30 & 0.01 & 0.02 & 4.21 & 0.03 & 0.05 & 4.15 \\
Cycles among mutual friends (\code{find\_rings}) & 0.98 & 0.95 & 0.05 & 0.89 & 0.96 & 0.29 & 0.82 & 0.94 & 0.50 & 0.99 & 0.99 & 0.03 & 0.83 & 0.99 & 0.47 \\
\quad + policy test, router scores (original) & 0.99 & 0.95 & 0.02 & 0.98 & 0.96 & 0.05 & 0.94 & 0.94 & 0.13 & 0.10 & 0.02 & 0.00 & 0.15 & 0.04 & 0.17 \\
\quad + policy test, fitted model (deployed) & 0.99 & 0.95 & 0.02 & 0.98 & 0.96 & 0.05 & 0.93 & 0.94 & 0.16 & 1.00 & 0.79 & 0.00 & 0.91 & 0.85 & 0.20 \\
\bottomrule
\end{tabular}

\vspace{2pt}
\parbox{0.96\textwidth}{\scriptsize Recall is over the 26 planted ring members; FPR is
over the 1,174 other students. Precision is counted as 0 in a trial that flags nobody;
``n/a'' marks detectors that flagged nobody in almost every trial. Structure is the same
in every regime, so the clustering row repeats. $^{\dagger}$Scored as a ranking cut at
the true number of ring members, which no deployment would know. Half-widths of 95\%
confidence intervals are at most 0.06 for precision or recall in the last three rows
and at most 0.12 elsewhere; Table~\ref{tab:summary} gives them for the deployed
detector.}
\end{table*}

\emph{Single signals see nothing.} Ring members, like everyone, have friends outside
their ring, so their closure stays under the 0.55 knee; and because farming errands
are cheap and members keep running ordinary errands, circulation stays under 0.6. The
clustering coefficient flags about five students per trial, mostly honest group
members, and almost never a ring member.

\emph{Graph-mining baselines find the marketplace, not the ring.} FRAUDAR's own
output, its densest block, was the whole campus (1,199 or 1,200 accounts) in every
trial. A ring of three to eight students who each run a dozen cheap errands for one
another is not dense compared with the core of an active marketplace, where everyone
pays everyone a little. Cut at the true number of ring members, the peeling order
still reached an F1 score of only 0.28--0.48. Label-propagation communities were about
the size of a hostel block, so no community crossed the circulation knee in any trial,
and the ranked version mostly flagged a whole block of honest students. Both methods
look for large dense or tightly-knit groups. The groups we are after are small, and
dense only in money that flows around a cycle.

\emph{Cycles need the friendship test.} Payment cycles without the friendship test
found almost every ring member but flagged 13 to 44 honest students per trial,
because repeated errands between nearby strangers also close cycles. Requiring every
leg to join mutual friends is what makes cycle evidence usable: \code{find\_rings}
reached an F1 score of $0.96\pm0.02$ in the routine regime and $0.99\pm0.01$ against
coordinated rings, which do not change the money pattern.

\emph{The router's false positives.} The ring search's false positives come from the
two cases it cannot see through: router-induced cycles (3.4 honest students per trial
in the busy regime) and honest favours (5.9 in the favours regime, 5.5 in the hard
one). Both versions of the policy test removed most of them. The fitted test removed
82\% in the busy regime and 68\% with honest favours, lifting F1 from 0.92 to 0.96 and
from 0.87 to 0.93, and lost no ring member in either.

\emph{Where the original test failed.} Against coordinated rings the router-score
test cleared 146 of the 149 cycles that contained ring members, and recall fell from
0.99 to 0.02. The fitted test cleared 37 of them, for a recall of $0.79\pm0.06$ at a
precision of 1.00. With honest favours added (the hard regime) the original test kept
a recall of 0.04 and the fitted test 0.85, at a precision of 0.91.
Section~\ref{sec:eval-policy} examines why.

\emph{Variability.} Recall below 1 in the routine, busy and favours regimes came from
seeds in which honest runners won enough of a ring's cheap errands to push one leg
under the qualification floor (on average 1.3 of 26 members missed in the routine
regime). Fig.~\ref{fig:sweeps}(b) shows how that floor depends on farming intensity.
The number of cycles raised per trial varied from four to nine, because some seeds
plant honest groups where the router closes a cycle.
